\documentclass{article}
\usepackage{amsmath,amssymb,amsthm}
\newtheorem{theorem}{Theorem}
\newtheorem{definition}{Definition}
\newtheorem{lemma}{Lemma}
\newtheorem{corollary}{Corollary}
\title{Linear Algebra II Home Assignment 1}
\author{
\begin{tabular}{c}
Anurag Bhunia (bmat2508@isibang.ac.in)\\[0.6em]
\end{tabular}
}
\date{}
\renewcommand{\thesubsection}{\arabic{subsection}}
\newcommand{\Col}{\operatorname{Col}}
\newcommand{\Null}{\operatorname{Null}}
\begin{document}
\maketitle
\begin{enumerate}
\item Consider a permutation $\sigma :=(2,5)(1,4,3)$ of $\{1,2,3,4,5,6\}$. Write down the matrices $P^{\sigma},P^{\sigma^{-1}}$. Verify that $\det{P^{\sigma}}=\epsilon(\sigma)$ for this permutation.
\end{enumerate}
\begin{proof}[Sol.]
Due to the way this permutation has been defined, it must swap the second and fifth coordinates of a vector in $\mathbb{R}^6$, and must take the first coordinate to the fourth, the fourth to the third and the third back to the first; all while keeping the sixth coordinate in its place.
Let the matrix for $\sigma$ be $P^{\sigma}$. Then the action of $P^{\sigma}$ on the vector 
$\begin{pmatrix}
1\\
2\\
3\\
4\\
5\\
6
\end{pmatrix}$ must be as follows:
$$P^{\sigma}
\begin{pmatrix}
1\\
2\\
3\\
4\\
5\\
6
\end{pmatrix}=
\begin{pmatrix}
3\\
5\\
4\\
1\\
2\\
6
\end{pmatrix}
$$
In that case $P^{\sigma}$ acts on the rows of a vector $\mathbf{v}\in\mathbb{R}^6$. We observe that the matrix 
$$
\begin{pmatrix}
0&0&1&0&0&0\\
0&0&0&0&1&0\\
0&0&0&1&0&0\\
1&0&0&0&0&0\\
0&1&0&0&0&0\\
0&0&0&0&0&1
\end{pmatrix}
$$
does just that, and since matrices for linear transformations are unique, we conclude that $P^{\sigma}$ is the matrix as described above.\\
Next we note that if a permutation matrix takes the $i$-th coordinate of a vector to the $j$-th coordinate then the inverse should take the $j$-th coordinate to the $i$-th one.
With this, we have that 
$$
P^{\sigma^{-1}}=
\begin{pmatrix}
0&0&0&1&0&0\\
0&0&0&0&1&0\\
1&0&0&0&0&0\\
0&0&1&0&0&0\\
0&1&0&0&0&0\\
0&0&0&0&0&1
\end{pmatrix}
$$\\
We can bring the permutation matrix back to the identity by 
\renewcommand{\qedsymbol}{}
\end{proof}
\end{document}